% =====================================================================
\section{Full Verification Prompt (\Acond{1})}
\label{app:prompt}
% =====================================================================
% Source: src/prompts/ablation_verification_prompts.py; the stored file
% outputs/verification_ablation_5747/A1_overlay_only/prompts/sample_000001.txt
% TODO: replace this transcription with an \input or verbatiminput of the
% stored prompt file so that it is byte-identical (copy it into paper/appendix/).
The prompt below is the text built for \Acond{1}. Other conditions replace the
``Input image'' block and the condition-specific block at the end
(\cref{tab:conditions}).

\begin{footnotesize}
\begin{verbatim}
You are an expert reviewer for aerial wild palm detection in UAV orthomosaic
imagery.

Your task is NOT to detect new palms.

Your task is to verify whether the highlighted YOLO detection contains a valid
wild palm.

--------------------------------------------------

Important instructions

The highlighted green bounding box indicates the candidate detection.

Focus primarily on the highlighted object.

Use surrounding context only if it provides useful evidence.

The YOLO confidence score is auxiliary information only.

Do NOT base your decision primarily on the confidence score.

Use visual evidence as the primary basis for your judgment.

--------------------------------------------------

Wild palm characteristics

Wild palms in aerial orthomosaic imagery often exhibit:

* a visible central crown
* radial or circular crown structure
* long feather-like fronds extending outward
* repeated symmetric leaf patterns
* texture distinguishable from surrounding broadleaf canopy
* fronds radiating from a central point

A palm may be partially occluded.

A partially visible palm can still be considered Reliable if recognizable palm
morphology is present.

Dense vegetation without recognizable palm morphology should NOT be considered
a valid palm.

--------------------------------------------------

Input image
The input is an orthomosaic patch where the background outside the target
region is slightly dimmed and a green bounding box highlights the candidate
detection.

--------------------------------------------------

Decision definitions

Reliable
The highlighted bounding box clearly contains a valid wild palm.

Uncertain
Evidence is ambiguous due to occlusion, overlap, low image quality, partial
visibility, or confusing vegetation.

Unreliable
The highlighted object is not a palm or the detection is clearly incorrect.

--------------------------------------------------

Reasoning requirements

Base your reasoning primarily on:

* crown structure
* frond pattern
* texture
* morphology
* spatial arrangement

Do not rely mainly on confidence score.

--------------------------------------------------

Judge only from the highlighted bounding box in the image.
Do not use YOLO confidence or bounding-box geometry.

--------------------------------------------------

Ablation condition: A1_overlay_only
Sample ID: sample_000001

Return ONLY valid JSON.
{
  "decision": "Reliable | Uncertain | Unreliable",
  "confidence_reasoning": "",
  "visual_reasoning": ""
}
\end{verbatim}
\end{footnotesize}
\noindent Note: bullet characters (\textbullet) in the source are shown as
\texttt{*} here, and long lines are wrapped for the page.
\VERIFY{Diff against the stored prompt file before submission.}

% =====================================================================
\section{Condition-Specific Prompt Blocks}
\label{app:prompt_blocks}
% =====================================================================
\TODO{Reproduce the ``Input image'' descriptions for \Acond{4} and \Acond{5} and the metadata blocks/instructions for \Acond{2}--\Acond{5} from \texttt{src/prompts/ablation\_verification\_prompts.py}.}

% =====================================================================
\section{Reproducibility Details}
\label{app:repro}
% =====================================================================
\TODO{Per-model table: checkpoint, upstream repository and revision (if recorded), dtype, attention backend, Python / \texttt{transformers} / \texttt{torch} versions, GPU, experiment ID, Slurm job IDs. Source: \texttt{docs/EXPERIMENT\_STATUS\_CANONICAL.md} §3, \texttt{docs/MODEL\_EXPANSION\_CAMPAIGN.md} ledger, \texttt{logs/campaign/submissions.csv}, per-run \texttt{\_run\_metadata/}.}

\TODO{Git commit(s) of the code used for each run and for each analysis script; Protocol-v2 \texttt{PROTOCOL.json} records the base commit of the re-score.}

% =====================================================================
\section{Qualification History}
\label{app:qualification}
% =====================================================================
\TODO{Summarize per-model qualification path and documented deviations (\texttt{docs/EXPERIMENT\_STATUS\_CANONICAL.md} §5). State that historical balanced-100 gates were alignment-based.}

% =====================================================================
\section{Additional Results}
\label{app:additional}
% =====================================================================
% TODO: AUTO-GENERATE FROM RESULTS
\TODO{Full per-model confusion matrices; full pairwise transition tables; A1@1{,}000 results for LLaVA-OneVision and Gemma~3 (separately; prior 0.928, not comparable with full-cohort numbers); Protocol~v1$\rightarrow$v2 comparison (provenance).}

% =====================================================================
\section{Semantic Review Interface and Criteria}
\label{app:semantic}
% =====================================================================
\TODO{Screenshot of the blind review app (\texttt{scripts/semantic\_review\_app.py}); labeling criteria; provenance hashes from \texttt{docs/SEMANTIC\_VALIDITY\_AUDIT.md} §G.}
